\documentclass[10pt,a4paper]{amsart}
\usepackage[margin=19mm]{geometry}
\usepackage{amsmath,amssymb,mathtools}
\usepackage[scr=rsfs,cal=euler]{mathalfa}
\usepackage{xfrac}
\usepackage[unicode,hidelinks,bookmarks=false]{hyperref}
\newcommand{\ie}{i.e.\ }
\newcommand{\cf}{cf.\ }
\newcommand{\resp}{respectively\ }
\newcommand{\Subsection}{\subsection}
\providecommand{\nobreakdash}{\nobreak\hbox{-}}
\setlength{\emergencystretch}{2em}
\multlinegap0pt
\usepackage{bm}
\usepackage{bbm}
\usepackage{dsfont}
\usepackage{xspace}
\usepackage{cancel}
\usepackage{mathtools}
\usepackage{amsthm}
\makeatletter
\@ifundefined{theorem}{\newtheorem{theorem}{Theorem}}{}
\@ifundefined{lemma}{\newtheorem{lemma}[theorem]{Lemma}}{}
\makeatother
\newcommand{\Ind}{\operatorname{Ind}}
\newcommand{\SYT}{{\operatorname{SYT}}}
\newcommand{\al}{{\alpha}}
\newcommand{\id}{{\operatorname{id}}}
\newcommand{\la}{{\lambda}}
\newcommand{\todo}[1]{\vspace{5 mm}\par \noindent
	\marginpar{\textsc{ToDo}} \framebox{\begin{minipage}[c]{0.95
				\textwidth}
			#1 \end{minipage}}\vspace{5 mm}\par}
\renewcommand\S{{\mathcal S}}
\title[Asymptotics of higher Lie characters]{Asymptotics of higher Lie characters}
\author{Ron M. Adin, Yuval Roichman, Natalia Tsilevich}
\date{Source: arXiv:2509.12904v1 (CC BY 4.0)}
\begin{document}
\maketitle

\noindent\emph{Source and licence.} Asymptotics of higher Lie characters. Version arXiv:2509.12904v1 (CC BY 4.0), distributed under the Creative Commons Attribution 4.0 International licence, as recorded in the arXiv licence metadata for this exact version and shown on the abstract page https://arxiv.org/abs/2509.12904v1. Full rights notice: licenses/arxiv-2509.12904-CC-BY-4.0.txt.\par\smallskip

\noindent\textit{Editorial scope.} Counted unit: the Theorem whose source label is \texttt{thm:blocks} in the article (source0.tex lines 3599--3609, source label \texttt{thm:blocks}), delivered together with the article's own complete original proof of it (source0.tex lines 3612--3773, one \texttt{proof} environment of 162 lines). No mathematics is written, repaired or re-derived: statement and proof are the article's own text line for line. Required context retained uncounted: eq:hlcformula (lines 676--679); thm:rectangular (lines 1137--1154). Every definition, display and label of those lines is retained unchanged. Omitted from this package and not redistributed: 1--675, 680--1136, 1155--3598, 3610--3611, 3774--4283. Nothing inside a retained excerpt is omitted or altered: every retained body is delivered exactly as the article's lines are written, so no omission receipt of any kind accompanies this package. Citations: the retained excerpts carry no citation command at all, so no bibliography entry of the article is transcribed and no reference is redistributed. Reference adaptation: none was needed and none was made; every reference inside the delivered unit resolves inside it, so no unresolved reference or citation is produced. Printed slips kept literal: no printed word, symbol, hypothesis or number of the counted statement or of its proof was altered. Layout and class: the article's own class is \texttt{amsart} with packages including fancyhdr, url, amsmath, amsthm, amssymb, tikz, hyperref, amsfonts, geometry; the wrapper is the lane's pinned \texttt{amsart} editorial wrapper at 10pt on A4, which restates the theorem-like environments the retained text needs and adds the article's own macro definitions that the retained text needs (\texttt{\textbackslash Ind}, \texttt{\textbackslash S}, \texttt{\textbackslash SYT}, \texttt{\textbackslash al}, \texttt{\textbackslash id}, \texttt{\textbackslash la}, \texttt{\textbackslash todo}), with extra wrapper packages (bm, bbm, dsfont, xspace, cancel, mathtools). The delivered numbering is the unit's own. Assets: no retained span contains a figure environment, a graphics-inclusion command, a PDF-page inclusion, a table or a TikZ picture; no image file is copied into this package, the package declares no asset dependency and no image file is redistributed with it. Licence: version arXiv:2509.12904v1 of this article is distributed under the Creative Commons Attribution 4.0 International licence, as recorded in the arXiv licence metadata for this exact version and shown on the abstract page https://arxiv.org/abs/2509.12904v1. A full rights notice is delivered at licenses/arxiv-2509.12904-CC-BY-4.0.txt. Characters: every retained excerpt is pure ASCII, so the delivered engine, XeLaTeX with its default Latin Modern text font, reports no missing character.
\begin{equation}\label{eq:hlcformula}
\psi^{\la}=\Ind_{\S_{m_1}[\S_1]\times\S_{m_2}[\S_2]\times\ldots}^{\S_n}
\left(\id_{m_1}[\psi^{(1)}]\otimes\id_{m_2}[\psi^{(2)}]\otimes\ldots\right),
\end{equation}

\begin{theorem}\label{thm:rectangular}
    Any sequence of higher Lie characters indexed by rectangular diagrams $\la=(m^k)$ with $m\to\infty$ and $k = o(m)$ tends to be regular. Furthermore,
    \begin{equation*}%\label{rectangular}
        \langle L_{(m^k)},s_\nu \rangle
        \sim \frac{f^\nu}{m^kk!} 
        \quad \text{as} \quad m \to \infty
    \end{equation*}
    for balanced (and hence for Plancherel-typical) $\nu \vdash mk$. 
    
%    More precisely, fix $C, c, \delta >0$. For any $\al > 0$ there exists $M > 0$, depending only on $\alpha$, $C$, $c$ and $\delta$,
%    such that
%    \begin{equation}\label{rectangularerror}
%        \langle L_{(m^k)},s_\nu \rangle
%        = \frac{f^\nu}{m^kk!} (1+R_\nu)
%        \quad \text{ with } \quad |R_\nu|\le m^{-\al}
%    \end{equation}
%    for any $m \ge M$, $k \le c m^{1 - \delta}$ and $\nu \vdash km$ satisfying $\nu_1, \nu'_1 \le C \sqrt{km}$.
\end{theorem}

\begin{theorem}\label{thm:blocks}
Let $\la=(m_1^{k_1} \cdots m_\ell^{k_\ell}) \vdash n=\sum k_im_i$, where $m_1 > \dots > m_\ell$ are distinct and
{\rm(a)} $k_i=o(m_i)$; 
{\rm(b)} $k_im_i\asymp n$; 
{\rm(c)} $\ell = o(n)$. 
Then this sequence of diagrams is regular:
$$
    \langle L_{\la},s_\nu \rangle \sim \frac{f^\nu}{z_\la} \qquad \text{as }n\to \infty
$$ 
for  typical  $\nu\vdash n$.
\end{theorem}

\begin{proof}
It follows from \eqref{eq:hlcformula} that
$$
L_\la=\prod_{i=1}^\ell L_{(m_i^{k_i})}.
$$
We will use the following formula: for arbitrary symmetric functions $f,g$,
\begin{equation}\label{eq:product}
    \langle fg,s_\nu \rangle
    = \sum_{\la\subset\nu} \langle s_\la, f \rangle \langle s_{\la/\mu}, g \rangle 
    =\sum_{\la,\mu\subset\nu} c^\nu_{\la\mu} \langle s_\la, f \rangle \langle s_\mu, g \rangle,
\end{equation}
and its obvious generalization to any number of factors.



Then
\[
    \langle L_\la,s_\nu \rangle
    = \langle \prod_{i=1}^\ell L_{(m_i^{k_i})}, s_\nu \rangle
    = \sum_{\nu_i \vdash k_im_i} c^\nu_{\nu_1,\ldots,\nu_\ell} \prod_{i=1}^\ell \langle L_{(m_i^{k_i})},s_{\nu_i} \rangle,
\]
    where $c_{\nu_1,\ldots,\nu_r}^\nu$ is a ``multi-Littlewood--Richardson coefficient'' equal to the coefficient of $s_\nu$ in the product
    $s_{\nu_1} \cdots s_{\nu_r}$. For typical  $\nu\vdash n$, the rows and columns of $\nu$ are $O(\sqrt n)=O(\sqrt{k_im_i})$ in view of assumption~(b). Since $c_{\nu_1,\ldots,\nu_r}^\nu\ne0$ only if $\nu_j\subset \nu$, the rows and columns of $\nu_i\vdash k_im_i$ are also $O(\sqrt{k_im_i})$.

Then, by Theorem~\ref{thm:rectangular}, for any $\al>0$
$$
\langle L_{(m_i^{k_i})},s_{\nu_i}\rangle=\frac{f^{\nu_i}}{m_i^{k_i}k_i!}(1+R_{\nu_i})\qquad\text{where}\qquad
|R_{\nu_i}|\le m_i^{-\al}
$$
for sufficiently large $m_i$. It follows from assumptions~(a) and~(b) that 
$k_i=o(\sqrt {n})$ and that if $n$ is large enough, then {\em all the $m_i$'s are large enough}.

\todo{NT:
The claim that all the $m_i$'s are large enough is false.

It suffices to assume $\ell = O(\ln n)$.
}

\todo{YR: 
I believe that we do not need all $m_i$-s to be proportional to $n$,  
and that the gap could be resolved. 

%If I understand correctly, for all parts but $\log \log n$ parts,  $m_i \ge \log n$. Letting $\alpha>1$ the sum of the errors on these parts tend to zero (since there are $\sim \log n$ parts).
 
%For the smallest parts, I think that their sizes grow like "a geometric sequence", 
%thus the sum of the errors is bounded. Furthermore, summing the errors over all but finite smallest parts, the sum could be close to 1 as we wish. 
%Is this correct?\\

%\smallskip

Explanation: I believe that the following lemma holds.

\begin{lemma}
for a random permutation, 
the expected size of the $i$-th cycle (when the cycles are ordered by length) is $m_i \sim n  e^{-i}$, as long as this size tends to infinity. 
\end{lemma}

Here is a heuristic:  The expected number of $k$-cycles  is $1/k$, 
hence 
%for any $m$, 
the expected number of cycles of length %$m<k\le e^i m$ is $i$. 
$e^{-i}n\le k\le n$ is $\sim \int_{ne^{-i}}^n {1\over x} dx= i$. 

Another heuristic: Construct the random permutation by choosing uniformly a letter $j$, then $\pi(j)$ then $\pi^2(j)$, etc., one can verify that the expected length of the resulting cycle is $\sim n/e$; continue recursively.

If the lemma is correct, then the sum of the errors is 
\[
\sum m_i^{-\alpha} \sim n^{-\alpha} \sum_{i=1}^{\log n} e^{i\alpha}  \sim 
n^{-\alpha} \int_{1}^{\log n} e^{\alpha x} dx \sim 
\alpha^{-1}. 
\]
This argument should be cleaned up by, 
first, taking in account in the lemma concentration, not just expectation (for this [Arratia-Tavare 92] could be helpful; see 
see also John Baez homepage, random permutations - Part 7). %https://math.ucr.edu/home/baez/permutations/permutations_7.html). 
Second,  in the displayed sum of errors,  
one should take in account 
the part sizes which are bounded.  

}

Thus, 
$m_i^{-1}\asymp k_in^{-1}=o(n^{-1/2})$.
Therefore, 
$$
\prod_{i=1}^\ell(1+R_{\nu_i})-1=\prod_{i=1}^\ell(1+o(n^{-\al/2}))-1=O(\ell n^{-\al/2})
=o(n^{-\al/2-1})
$$
and
$$
\prod\langle L_{(m_i^{k_i})},s_{\nu_i}\rangle=\frac{1}{\prod m_i^{k_i}k_i!}\left(1+o(n^{-\al/2-1})\right)\prod f^{\nu_i}.
$$
Then 
$$
\langle L_{\la},s_\nu\rangle=\frac{1}{\prod m_i^{k_i}k_i!}\left(1+o(n^{-\al/2-1})\right)
\sum_{\nu_i\vdash k_im_i}c^\nu_{\nu_1,\ldots\nu_\ell}\prod f^{\nu_i}.
$$
But
$$
\sum_{\nu_i\vdash k_im_i}c^\nu_{\nu_1,\ldots\nu_\ell}\prod f^{\nu_i}=\sum_{\nu_i\vdash k_im_i}c^\nu_{\nu_1,\ldots\nu_\ell}\prod\langle h_1^{k_im_i},s_{\nu_i}\rangle
=\langle h_1^{\sum k_im_i},s_\nu\rangle=f^\nu,
$$
and the theorem follows.

\todo{YR: Another suggestion for a more refined count, which may fix the gap: 

First, fix $0<\epsilon << 1$. A partition $\nu\vdash n$ is called {\em 
$\epsilon$-good} if 
 with $\nu_1+\nu'_1\le \epsilon n$. Otherwise it is {\em $\epsilon$-bad}.

Notice that for an $\epsilon$-good partition $\nu\vdash n$,  $f^\nu \ge c^n$ for sufficiently large $n$ and constant $c$ which depends on $\epsilon$. 

Hence, by Theorem 3.6 and proof of Proposition 4.1, the error term for such partitions is exponentially decreasing with $n$. 

Next (assume at the moment that all in $\lambda$ parts are distinct)
\[
\langle L_{\la},s_\nu\rangle=\frac{1}{\prod m_i}
\sum_{\nu_i\vdash m_i}c^\nu_{\nu_1,\ldots\nu_\ell}\prod \left(1+Error(\nu_i)\right)f^{\nu_i}.
\]
Thus, for each series of partitions in the sum $\nu_1\vdash m_1,\dots,\nu_\ell\vdash m_\ell$
\[
\prod \left(1+Error(\nu_i)\right)f^{\nu_i}=\prod_i \left(1+c^{-m_i}\right)f^{\nu_i} \prod_{\nu_i\text{ is $\epsilon$-bad}} \left(1+m_i^{-1}\right)f^{\nu_i}.  
\]
It follows that 
\[
\sum_{\nu_i\vdash m_i}c^\nu_{\nu_1,\ldots\nu_\ell}\prod \left(1+Error(\nu_i)\right)f^{\nu_i}
\]
\begin{equation}\label{eq:1111}
\le\sum_{\forall i\, \nu_i\vdash m_i\atop \nu_i \text{ is $\epsilon$-good}}c^\nu_{\nu_1,\ldots\nu_\ell}\prod (1+c^{-m_i}) f^{\nu_i} + \sum_{\exists i\, \nu_i\vdash m_i\atop \nu_i\text{ is $\epsilon$-bad} 
}c^\nu_{\nu_1,\ldots\nu_\ell}\prod (1+m_i^{-1}) f^{\nu_i}. 
\end{equation}
Recall that
\[
\sum_{\nu_i\vdash m_i}c^\nu_{\nu_1,\ldots\nu_\ell}\prod f^{\nu_i}=f^\nu,
\]
Hence, the first sum in \eqref{eq:1111} is  $\le (1+\sum_i c^{- m_i})f^\nu$. 

For the second sum, recall that the summand $c^\nu_{\nu_1,\nu_2} f^{\nu_1}f^{\nu_2}$ is equal to the number of SYT in $\SYT(\nu)$, for which the letters $\le |\nu_1|$ form a sub-tableau of shape $\nu_1$, and the rest of the letters, after operating jdt, form a tableaux of shape $\nu_2$.
Similar phenomenon holds for all series $\nu_1\vdash m_1,\dots,\nu_\ell\vdash m_\ell$. 




The following Observation completes the proof of the upper bound. 

{\bf Observation.} The probability  that the first $\nu_1$ letters in a  SYT of a typical shape shape $\nu$ is  $\epsilon$-bad is exponentillay decreasing with $m_1$.   
By symmetry of LR coefficients, this holds for every $1\le i\le \ell$. 

{\bf Proof of Observation.} First, observe that the shape $\nu_1$ formed by the prefix of 
a SYT of a typical shape $n$ is also typical as long as $m_1=|\nu_1|$ tends to infinity. Now, 
the number of SYT of a typical shape $\nu_1$ is asymptotically  
equal to 
\[
\frac{\sqrt {m_1!}}{e^{\sqrt m_1}}.
\]
The number of SYT of shape $\nu_1$ that are $\epsilon$-bad is 
\[
\le \binom{m_1}{\epsilon m_1} \binom{\epsilon m_1}{\epsilon m_1/2} {\sqrt {(1-\epsilon)m_1!}\over e^{\sqrt{(1-\epsilon m_1}} }. 
\]
The quotient is exponentially decreasing with $m_1$. One concludes that the second sum\\ $\le  \cdot f^\nu \prod\limits_i e^{\sqrt m_i}\gamma^{-m_i}(1+m_i^{-1})$. The proof of lower bound is similar. 
}

\end{proof}

\end{document}
